% ===============================
% SEZIONE 1: Fasci di rette
% ===============================

\begin{question}[ID=1.1, subtitle={Fascio $(k+1)x-(k+2)y+2=0$: tipo e centro}, tags={fascio di rette}, difficulty=facile]
Dato il fascio di rette di equazione $(k+1) \ x - (k+2) \ y +2 = 0$, determinare le generatrici del fascio, il tipo di fascio ed eventualmente il suo centro.
\begin{solution}
Si raccoglie il parametro: $k(x-y)+(x-2y+2)=0$. Le \textbf{generatrici} sono $r_1:\ x-y=0$ e $r_2:\ x-2y+2=0$. Essendo non parallele, il fascio è \textbf{proprio} e il centro è la loro intersezione: da $x=y$ in $x-2y+2=0$ si ha $-y+2=0\Rightarrow y=2$, quindi il centro è $C(2,2)$.
\end{solution}
\end{question}

\begin{question}[ID=1.2, subtitle={Fascio $(k+1)x+(k-1)y+2k+3=0$: tipo e centro}, tags={fascio di rette}, difficulty=facile]
Dato il fascio di rette di equazione $(k+1) \ x + (k-1) \ y + 2k + 3 = 0$, determinare le generatrici del fascio, il tipo di fascio ed eventualmente il suo centro.
\begin{solution}
Raccogliendo $k$: $k(x+y+2)+(x-y+3)=0$. Le \textbf{generatrici} sono $r_1:\ x+y+2=0$ e $r_2:\ x-y+3=0$, non parallele: fascio \textbf{proprio}. Centro: sommando le due, $2x+5=0\Rightarrow x=-\tfrac52$; sottraendo, $2y-1=0\Rightarrow y=\tfrac12$. Centro $C\!\left(-\tfrac52,\tfrac12\right)$.
\end{solution}
\end{question}

\begin{question}[ID=1.3, subtitle={Fascio $(k+1)x-(k+2)y+2=0$: determinazione di $k$}, tags={fascio di rette}, difficulty=facile]
Dato il fascio di rette di equazione $(k+1) \ x - (k+2) \ y +2 = 0$, determinare il parametro $k$ affinché la retta corrispondente del fascio:
\begin{enumerate}
    \item passi per il punto $P(-2,4)$
    \item sia parallela all'asse $y$
    \item sia parallela alla retta di equazione $4x+2y-3=0$
    \item sia perpendicolare alla retta di equazione $x-5y-3=0$
    \item intersechi l'asse $x$ in un punto di ascissa negativa e l'asse $y$ in un punto di ordinata positiva.
\end{enumerate}
\begin{solution}
Coefficiente angolare della generica retta: $m=\dfrac{k+1}{k+2}$.
\begin{enumerate}
\item $P(-2,4)$: $-2(k+1)-4(k+2)+2=0\Rightarrow-6k-8=0\Rightarrow k=-\tfrac43$.
\item Parallela all'asse $y$: coefficiente di $y$ nullo, $-(k+2)=0\Rightarrow k=-2$ (retta $x=2$).
\item Parallela a $4x+2y-3=0$ (pendenza $-2$): $\dfrac{k+1}{k+2}=-2\Rightarrow k=-\tfrac53$.
\item Perpendicolare a $x-5y-3=0$ (pendenza $\tfrac15$, perpendicolare $-5$): $\dfrac{k+1}{k+2}=-5\Rightarrow k=-\tfrac{11}{6}$.
\item Intercetta $x$: $x=-\dfrac{2}{k+1}<0\Rightarrow k>-1$; intercetta $y$: $y=\dfrac{2}{k+2}>0\Rightarrow k>-2$. Insieme: $k>-1$.
\end{enumerate}
\end{solution}
\end{question}

\begin{question}[ID=1.4, subtitle={Fascio $kx+(k-1)y+1=0$: determinazione di $k$}, tags={fascio di rette}, difficulty=facile]
Dato il fascio di rette di equazione $k\ x + (k-1) \ y +1 = 0$, determinare il parametro $k$ affinché la retta corrispondente del fascio:
\begin{enumerate}
    \item passi per il punto $P(2,0)$
    \item sia parallela all'asse $y$
    \item sia parallela alla bisettrice del primo e terzo quadrante
    \item sia perpendicolare alla retta di equazione $x-5y-3=0$
    \item formi con l'asse $x$ un angolo di $30^\circ$.
\end{enumerate}
\begin{solution}
Coefficiente angolare: $m=-\dfrac{k}{k-1}$.
\begin{enumerate}
\item $P(2,0)$: $2k+1=0\Rightarrow k=-\tfrac12$.
\item Parallela all'asse $y$: coefficiente di $y$ nullo, $k-1=0\Rightarrow k=1$ (retta $x=-1$).
\item Parallela alla bisettrice $y=x$ (pendenza $1$): $-\dfrac{k}{k-1}=1\Rightarrow k=\tfrac12$.
\item Perpendicolare a $x-5y-3=0$ (pendenza $\tfrac15$, perpendicolare $-5$): $-\dfrac{k}{k-1}=-5\Rightarrow k=\tfrac54$.
\item Angolo di $30^\circ$, $m=\pm\tfrac{1}{\sqrt3}$: da $-\dfrac{k}{k-1}=\tfrac{1}{\sqrt3}$ si ha $k=\dfrac{\sqrt3-1}{2}$; da $-\dfrac{k}{k-1}=-\tfrac{1}{\sqrt3}$ si ha $k=-\dfrac{\sqrt3+1}{2}$.
\end{enumerate}
\end{solution}
\end{question}

% ===============================
% SEZIONE 2: Problemi
% ===============================
